\documentclass[10pt,a4paper]{amsart}
\usepackage[margin=19mm]{geometry}
\usepackage{amsmath,amssymb,mathtools}
\usepackage[scr=rsfs,cal=euler]{mathalfa}
\usepackage{xfrac}
\usepackage[unicode,hidelinks,bookmarks=false]{hyperref}
\newcommand{\ie}{i.e.\ }
\newcommand{\cf}{cf.\ }
\newcommand{\resp}{respectively\ }
\newcommand{\Subsection}{\subsection}
\providecommand{\nobreakdash}{\nobreak\hbox{-}}
\setlength{\emergencystretch}{2em}
\multlinegap0pt
\usepackage{bm}
\usepackage{bbm}
\usepackage{dsfont}
\usepackage{xspace}
\usepackage{cancel}
\usepackage{mathtools}
\usepackage{amsthm}
\makeatletter
\@ifundefined{equation}{\newtheorem{equation}{Equation}}{}
\@ifundefined{theorem}{\newtheorem{theorem}{Theorem}}{}
\@ifundefined{corollary}{\newtheorem{corollary}[equation]{Corollary}}{}
\@ifundefined{definition}{\newtheorem{definition}[equation]{Definition}}{}
\@ifundefined{lemma}{\newtheorem{lemma}[equation]{Lemma}}{}
\@ifundefined{proposition}{\newtheorem{proposition}[equation]{Proposition}}{}
\makeatother
\newcommand{\RR}{\mathds{R}}
\newcommand{\cC}{\mathcal{C}}
\newcommand{\cH}{\mathcal{H}}
\newcommand{\cL}{\mathcal{L}}
\newcommand{\cM}{\mathcal{M}}
\newcommand{\cN}{\mathcal{N}}
\newcommand{\cV}{\mathcal{V}}
\title[The Bounded Diameter Conjecture and Sharp Geometric Estimate]{The Bounded Diameter Conjecture and Sharp Geometric Estimates for Mean Curvature Flow}
\author{Yiqi Huang, Wenshuai Jiang}
\date{Source: arXiv:2510.17060v2 (CC BY 4.0)}
\begin{document}
\maketitle

\noindent\emph{Source and licence.} The Bounded Diameter Conjecture and Sharp Geometric Estimates for Mean Curvature Flow. Version arXiv:2510.17060v2 (CC BY 4.0), distributed under the Creative Commons Attribution 4.0 International licence, as recorded in the arXiv licence metadata for this exact version and shown on the abstract page https://arxiv.org/abs/2510.17060v2. Full rights notice: licenses/arxiv-2510.17060-CC-BY-4.0.txt.\par\smallskip

\noindent\textit{Editorial scope.} Counted unit: the Lemma whose source label is \texttt{l:c cover} in the article (source0.tex lines 964--973, source label \texttt{l:c cover}), delivered together with the article's own complete original proof of it (source0.tex lines 976--1034, one \texttt{proof} environment of 59 lines). No mathematics is written, repaired or re-derived: statement and proof are the article's own text line for line. Required context retained uncounted: l:near cylindrical also cylindrical (lines 540--542); p:almost cylindrical rigidity (lines 724--737); p:propagation-cylinder (lines 751--761); d:cylindricalregion (lines 800--809); l:cylindrical region structure (lines 818--823); c:measure upper bound (lines 850--853). Every definition, display and label of those lines is retained unchanged. Omitted from this package and not redistributed: 1--539, 543--723, 738--750, 762--799, 810--817, 824--849, 854--963, 974--975, 1035--1635. Nothing inside a retained excerpt is omitted or altered: every retained body is delivered exactly as the article's lines are written, so no omission receipt of any kind accompanies this package. Citations: the retained excerpts carry no citation command at all, so no bibliography entry of the article is transcribed and no reference is redistributed. Reference adaptation: none was needed and none was made; every reference inside the delivered unit resolves inside it, so no unresolved reference or citation is produced. Printed slips kept literal: no printed word, symbol, hypothesis or number of the counted statement or of its proof was altered. Layout and class: the article's own class is \texttt{amsart} with packages including times, amsmath, amsfonts, amstext, amssymb, amsbsy, amsopn, amsthm, eucal, txfonts, dsfont, graphicx, hyperref, accents; the wrapper is the lane's pinned \texttt{amsart} editorial wrapper at 10pt on A4, which restates the theorem-like environments the retained text needs and adds the article's own macro definitions that the retained text needs (\texttt{\textbackslash RR}, \texttt{\textbackslash cC}, \texttt{\textbackslash cH}, \texttt{\textbackslash cL}, \texttt{\textbackslash cM}, \texttt{\textbackslash cN}, \texttt{\textbackslash cV}), with extra wrapper packages (bm, bbm, dsfont, xspace, cancel, mathtools). The delivered numbering is the unit's own. Assets: no retained span contains a figure environment, a graphics-inclusion command, a PDF-page inclusion, a table or a TikZ picture; no image file is copied into this package, the package declares no asset dependency and no image file is redistributed with it. Licence: version arXiv:2510.17060v2 of this article is distributed under the Creative Commons Attribution 4.0 International licence, as recorded in the arXiv licence metadata for this exact version and shown on the abstract page https://arxiv.org/abs/2510.17060v2. A full rights notice is delivered at licenses/arxiv-2510.17060-CC-BY-4.0.txt. Characters: every retained excerpt is pure ASCII, so the delivered engine, XeLaTeX with its default Latin Modern text font, reports no missing character.
\begin{lemma}\label{l:near cylindrical also cylindrical}
    For any $\delta>0$, $\eta\le \eta_0(\delta)$, we have the following. Suppose $X$ is $(\eta,r)$-cylindrical with respect to $\cL_X = (x+L_X,t_X)$. Then any $Y \in Q_{\eta r}(\cL_X) \cap Q_{2r}(X)$ is $(\delta,r)$-cylindrical with respect to $(y+L_X, t_Y)$.  
\end{lemma}

\begin{proposition}\label{p:almost cylindrical rigidity}
    Let $\cM \subset \RR^3 \times I$ be a bounded almost regular flow. For any $\delta>0$, $r\le r_0(\cM,\delta)$ and $\eta\le \eta_0(\cM,\delta)$, the following holds for any $X_0$ with $t_{X_0} \ge 4\eta^{-2}r^2$. Suppose $ X = (x,t_X)\in \cV_{\eta,r}(X_0;\lambda)$ for some $\lambda \in [\lambda_1-\lambda_{gap},\lambda_1 +  \lambda_{gap}]$. Then we have the following
    \begin{enumerate}
        \item $X$ is $(\delta,s)$-cylindrical with respect to $\cL_{X} =(x+L_X,t_X)$ for any $s\in [\delta^2 r,r]$. Equivalently, in the normalized flow $\cM_{X,s'}$ for any $s'\in[\delta^3 r, \delta^{-1}r]$, the time-slice $\{t=-1\}$ inside the ball $B_{\delta^{-2}}(0)$ is a $C^{2,\alpha}$-graph with norm $\le\delta^2$ over a round cylinder with axis $L_X$.
        \item The $\lambda$-pinched set $\cV_{\eta,r}(X_0;\lambda)$ lies near the axis $\cL_{X}$:
        \begin{equation*}
            \cV_{\eta,r}(X_0;\lambda)\subset Q_{\delta r}(\cL_X)\cap Q_{2r}(X_0)\,.
        \end{equation*}
        \item For any $Y \in \cV_{\eta,r}(X;\lambda)$ that is $(\delta,r)$-cylindrical with respect to $\cL_Y=(y+L_Y,t_Y)$, we have
        \begin{equation*}
            d_H(\cL_X \cap Q_{2r}(X_0) , \cL_Y \cap Q_{2r}(X_0)) \le \delta^2 r ~~ \text{ and } ~~ d_H(L_X\cap B_1(0), L_Y \cap B_1(0)) \le \delta^2.
        \end{equation*}
    \end{enumerate}
\end{proposition}

\begin{proposition}\label{p:propagation-cylinder}
Let $\cM\subset\RR^3\times I$ be a bounded almost regular flow. Let $\delta>0$ and $0<r_1<r_2\le r_0(\cM,\delta)$. Then there exists $\eta_0=\eta_0(\cM,\delta)>0$ such that for any $0<\eta\le\eta_0$ the following holds. Suppose $X_0$ satisfies $t_{X_0}\ge4\eta^{-2}r_2^2$, and that there exists $X\in Q_{2r_2}(X_0)$ and some
\[
\lambda\in[\lambda_1-\lambda_{\mathrm{gap}},\lambda_1+\lambda_{\mathrm{gap}}]
\]
with
\[
\Theta_X(\eta r_1)\ \ge\ \lambda-\eta^2\ \ge\ \Theta_X(\eta^{-1}r_2)-\eta^2.
\]
Then $X$ is $(\delta,s)$-cylindrical with respect to some fixed axis $\cL_X$ for every $s\in[r_1,r_2]$.
\end{proposition}

\begin{definition}\label{d:cylindricalregion}
Let $\cM \subset \RR^3 \times I$ be a bounded almost regular flow. Let $\mathcal{C} \subset Q_{2r}(X_0)$ be a closed subset, and let $r_X : \mathcal{C} \rightarrow \RR_{\ge0}$ be a radius function and $\tau \le 10^{-10}$. The set $\mathcal{N} = Q_{2r}(X_0) \setminus \Bar{Q}_{r_X}(\cC)$ is called a $(\delta, r)$-cylindrical region if for any $X\in \cC$ the following hold:
\begin{enumerate}
    \item $\{Q_{\tau^2 r_X}(X) \}$ are pairwise disjoint for $X\in \cC$;
    \item for each $r_X \leq s \leq r $, $X$ is $
    (\delta,s)$-cylindrical with respect to $\mathcal{L}_{X} = L_X \times \{t_X\}$;
    \item for each $s \geq r_X$ with $Q_{2s}(X)\subset Q_{2r}(X_0)$, we have $\mathcal{L}_{X} \cap Q_s(X) \subset Q_{\tau s}(\mathcal{C})$ and $\mathcal{C} \cap Q_s(X) \subset Q_{\tau s}(\mathcal{L}_{X})$.
\end{enumerate}
Here $\Bar{Q}_{r_X}(\cC) \equiv \cup_{X\in \cC}\bar Q_{r_X}(X)$ is the union of the closure of central balls. 
\end{definition}

\begin{lemma}\label{l:cylindrical region structure}
    For any $X\in \cC$, the projection map $\pi_X : \cC \to \cL_{X}$ is a bi-Lipschitz map, i.e. 
\begin{equation}\label{e:bilipschitz of centers}
    (1-2\tau)d(X_1,X_2) \le    ||\pi_X(X_1) - \pi_X(X_2)|| \le d(X_1,X_2), \text{ for any } X_1, X_2 \in \cC.
\end{equation}
\end{lemma}

\begin{corollary}\label{c:measure upper bound}
Let $\cN \subset Q_{2r}(X_0) \setminus \Bar{Q}_{r_X}(\cC)$ be a $(\delta,r)$-cylindrical region. Then there exists some $C(\tau)$ such that 
    $\sum_{X\in \cC_+} r_X + \cH^1_P(\cC_0)  \le C(\tau) r$. 
\end{corollary}

\begin{lemma}[Covering of cylindrical balls]\label{l:c cover}
 For any $\delta \le \delta_0$, $\tau<10^{-10}$, $r \le r_0(\cM,\delta,\tau)$ and $\eta \le \eta(\cM,\delta,\tau)$, the following holds. Let $X_0$ with $t_{X_0}\ge 4\eta^2 r^{-2}$. Assume that $\sup_{Y\in Q_{2r}(X_0)}\Theta_Y(\eta^{-1}r) \le \lambda$ and that there exists some $X \in \cV_{\eta,r}(X_0;\lambda)$ for some $\lambda \in [\lambda_1 -\lambda_{gap}, \lambda_1 + \lambda_{gap}]$. Then we have the following decomposition
    \begin{equation*}
        Q_r(X_0) \subset (\cC_0 \cup \cN)  \cup \bigcup_{e \in E} Q_{r_e}(X_e),
    \end{equation*}
    where $\cN \subset Q_{2r}(X_0)$ is a $(\delta,r)$-cylindrical region and that any $Y \in Q_{2r_e}(X_e)$ satisfies that $\Theta_Y(\eta^{-1} r_e) < \lambda - \eta^2$. Moreover, $\cC_0$ is contained in an embedded Lipschitz submanifold of dimension 1 and there exists some constant $C(\tau,\eta)$ such that
    \begin{equation*}
        \sum_{e\in E} r_e +  \cH_P^1(\cC_0) \le C(\tau,\eta) r.
    \end{equation*}
\end{lemma}

\begin{proof}
    Fix some $\lambda \in [\lambda_1 -\lambda_{gap}, \lambda_1 + \lambda_{gap}]$. Since $X \in \cV_{\eta,r}(X_0;\lambda)$, by Proposition \ref{p:almost cylindrical rigidity}, $X$ is $(\eta',s)$-cylindrical with respect to $\cL_{X}= (x+L_X, t_X)$ for any $s \in [\tau^2 r,r]$, provided $\eta\le \eta(\eta')$ small enough. By Lemma \ref{l:near cylindrical also cylindrical}, any $Y \in Q_{\eta' s}(\cL_{X}) \cap Q_{2s}(X)$ is $(\delta, s)$-cylindrical with respect to $\cL_Y = (y+L_X, t_Y)$ for any $s \in [\tau^2 r, r]$, provided $\eta'\le \eta'(\delta,\tau)$ small enough.
    
    Consider a covering of $Q_{\eta' r}(\cL_X) \cap Q_{2r}(X)$ by balls $\{Q_{\tau r_1/10}(Y^1_i)\}_{i \in I_1}$ with $Y^1_i \in \cL_X$, $r_1=\tau r$ and $Q_{4\tau^2 r_1}(Y^1_i) \cap Q_{ 4\tau^2 r_1}(Y_j) = \emptyset$ whenever $Y^1_i \neq Y^1_j$. We define $\cC^1 = \cup_{i\in I_1} Y^1_i$ and then it is easy to see that $\cN^1 = Q_{2r}(X) \setminus \Bar{Q}_{r_1}(\cC^1)$ is a $(\delta,r)$-cylindrical region, provided $\eta$ small enough and the third condition in Definition \ref{d:cylindricalregion} satisfying with $\tau'\le \tau/2$. This is our first step covering. 
    
    Now we examine each ball $Q_{r_1}(Y^1_i)$. We classify the balls 
    \begin{equation*}
    \begin{split}
        E^1 &= \{ Y^1_i : \Theta_Y(\eta r_1) < \lambda - \eta^2 \text{ for any } Y \in Q_{2r_1}(Y^1_i) \}. \\
        F^1 &= \{Y^1_i\} \setminus E^1.
    \end{split}
    \end{equation*}
    
Then the ball $Q_{r_1}(Y_i^1)$ with center in $E^1$ has definite entropy drop for every point in $Q_{2r_1}(Y_i^1)$. We will see that these balls will be covered by $e$-ball as in the statement of the lemma. So we will not need to handle such balls until the final step. 
For each $Y^1_i \in F^1$, there exists some $X^2 \in Q_{r_1}(Y_i)$ such that $\Theta_{X^2}(\eta r_1) \ge \lambda -\eta^2$. By Proposition \ref{p:propagation-cylinder}, we have $X^2$ is $(\eta',s)$-cylindrical with respect to the fixed $\cL_{X^2} = (x^2 + L_{X^2}, t_{X^2})$ for any $s \in [\tau^2 r_1, r]$. By Lemma \ref{l:near cylindrical also cylindrical} and Proposition \ref{p:propagation-cylinder}, any $Y \in Q_{\eta' s}(\cL_{X^2})\cap Q_{2s}(Y^1_i)$ is $(\delta,s)$-cylindrical with respect to $\cL_Y = (y + L_{X^2}, t_Y)$ for any $s\in[\tau^2 r_1, r]$.
    
    
    Consider the covering of $Q_{\eta' r_1}(\cL_{X^2}) \cap Q_{2r_1}(Y^1_i)$ by balls $\{Q_{\tau r_2/10}(Y^2_i)\}$ with $Y^2_i \in \cL_{X^2}$, $r_2 = \tau^2 r$ and $Q_{4\tau^2 r_2}(Y^2_i) \cap Q_{4\tau^2 r_2}(Y^2_j) =\emptyset$ whenever $Y_i^2 \neq Y^2_j$. We do this covering for every $Y^1_i \in F^1$. Then we can collect all the centers and define $\cC^2 = E^1 \cup \bigcup_i \{Y^2_i\}$. We can check that $\cN^2 = Q_{2r}(X_0) \setminus \big(\bigcup_{e \in E^1}Q_{r_e}(X_e) \cup \bigcup_i \Bar{Q}_{r_2}(Y^2_i) \big)$ is a $(\delta,r)$-cylindrical region. Actually, the first two properties are satisfied by construction. We prove the third one for each $Y^2_i$ as the points in $E^1$ are checked in the last step. Note that $r_{Y^2_i} = r_2$. Since $\cL_{Y^2_i} = \cL_{X^2}$ inside the ball $Q_{2r_1}(Y^1_i)$, then we have 
$\text{ $\cL_{Y^2_i}\cap Q_s(Y^2_i) \subset Q_{\tau s/5}(\cC^2)$ and $\cC^2 \cap Q_s(Y^2_i) \subset Q_{\tau s/5}(\cL_{Y^2_i})$ }$
    for any $s\in [r_2,r_1]$. And for $s\in[r_1,r]$, by Proposition \ref{p:almost cylindrical rigidity}, we have $d_H(\cL_{Y_i^2}\cap Q_{s}(Y^2_i), \cL_{Y}\cap Q_s(Y^2_i)) \le \eta' s$ for any $Y \in \cC^2$ if $\eta\le \eta(\eta')$. We can pick $Y \in \cC^1$ and then the third property is checked using the third property in the cylindrical region $\cN^1$. Therefore, it follows that $\cN^2$ is a $(\delta, r)$-cylindrical region with 
\begin{align}\label{e:bettercovering5.8}
    \text{ $\cL_{Y}\cap Q_s(Y) \subset Q_{\tau s/5}(\cC^2)$ and $\cC^2 \cap Q_s(Y) \subset Q_{\tau s/5}(\cL_{Y})$ }
    \end{align}
    for any $Y\in \cC^2$ and $r\ge s\ge r_Y$. 

    Next we can classify the balls $Q_{r_2}(Y^2_i)$ and define
    \begin{equation*}
    \begin{split}
        E^2 &= E^1 \cup \{Y^2_i : \Theta_Y(\eta r_2) < \lambda - \eta^2 \text{ for any } Y \in Q_{2r_2}(Y^2_i) \} \\
        F^2 &= \{Y^2_i\} \setminus E^2.
    \end{split}
    \end{equation*}
    
We can iterate this process to obtain the covering
\begin{equation*}
    Q_r(X_0) \subset \cN^k \cup \bigcup_{e \in E^k} Q_{r_e}(X_e) \cup \bigcup_{Y^k_i \in F^k} Q_{r_k}(Y^k_i)
\end{equation*}
where 
\begin{equation*}
\begin{split}
    E^k &= E^{k-1} \cup \{Y^k_i: \Theta_Y(\eta r_k) < \lambda -\eta^2 \text{ for any } Y\in Q_{2r_k}(Y^k_i) \} \\
    F^k &= \{Y^k_i\} \setminus E^k
\end{split}
\end{equation*}
and $$\cN^k = Q_{2r}(X_0) \setminus \Big( \bigcup_{e\in E^k} \Bar{Q}_{r_e}(X_e) \cup \bigcup_{Y^k_i \in F^k} \Bar{Q}_{r_k}(Y^k_i) \Big)$$
is a $(\delta,r)$-cylindrical region with $\cC^k = E^k \cup F^k$. 

Now let $k \to \infty$ and thus $r_k \to 0$ and $F^k \to \cC_0$ in the Hausdorff metric. Moreover, for each $Y \in \cC_0$, we have $\lim_{s\to 0} \Theta_Y(s) \ge \lambda- \eta^2$. By Proposition \ref{p:propagation-cylinder}, $Y$ is $(\eta',s)$-cylindrical with respect to the fixed $\cL_Y = (y + \cL_Y, t_Y)$ for any $s\in(0, r]$. In particular, $Y$ is a singular point.  Hence we obtain the covering
\begin{equation*}
    Q_r(X_0) \subset \cN^k \cup \cC_0 \cup \bigcup_{e\in E}Q_{r_e}(X_e),
\end{equation*}
where $E = \cup_k E^k$ and $\cN = Q_{2r}(X_0) \setminus \big( \bigcup_{e\in E}Q_{r_e}(X_e) \cup \cC_0)$. We check that $\cN$ is a $(\delta,r)$-cylindrical region with $\cC = E \cup \cC_0$. Again, the first two properties are satisfied by construction. We check the third one. If $Y\in E^k$, then the third property is satisfied using the fact that $\cN^k$ is a $(\delta,r)$-cylindrical region and Proposition \ref{p:almost cylindrical rigidity}. We check the point $Y \in \cC_0$. Take any $r_k>r_Y=0$. Then there exists some $Y^k_i \in \cC^k$ such that $d(Y^k_i,Y) \le r_k$. Since $\cN^k$ is a $(\delta,r)$-cylindrical region satisfying \eqref{e:bettercovering5.8} for $\cC^k$, we have $\cL_{Y^k_i}\cap Q_{2r_k}(Y^k_i) \subset Q_{\tau r_k/5}(\cC^k)$ and $\cC^k \cap Q_{2r_k}(Y^k_i) \subset Q_{\tau r_k/5}(\cL_{Y^k_i})$. By proposition \ref{p:almost cylindrical rigidity}, we have $\cL_{Y^k_i}$ is $\eta'r_k$-close to $\cL_Y$ if $\eta\le\eta(\eta')$. This proves the third property for $\cL_{Y}$ at the scale $r_k$.

By  Lemma \ref{l:cylindrical region structure} and Corollary \ref{c:measure upper bound}, we obtain the Lipschitz structure of $\cC_0$ and the $1$-content estimate. To complete the proof, we still need to recover each ball $Q_{2r_e}(X_e)$ using $\{Q_{\eta^2 r_e}(X_e^j)\}_{1 \le j \le C_0}$ with $Q_{\eta^2 r_e/10}(X_e^{j_1}) \cap Q_{\eta^2 r_e /10}(X_e^{j_2}) = \emptyset$ provided $j_1 \neq j_2$. We denote the set $\tilde{E}$ to be the set of all those center $\{X_e^j\}_{e\in E; 1\le j \le C_0}$. And we define the new radius $\tilde{r}_e^j = \eta^2 r_e$. This implies that $\Theta_Y(\eta^{-1} \tilde{r}_e) < \lambda - \eta^2$ for any $Y \in Q_{2 \tilde{r}^j_e}(X^j_e)$. And moreover we have the estimate
\begin{equation*}
    \sum_{e\in E} \sum_{j=1}^{C_0}  r_e^j \le C(\tau,\eta)r.
\end{equation*}
Hence we complete the whole proof.
\end{proof}

\end{document}
